\section{Using the collision matrices to simulate linear relaxation}
\label{qt:dynamics}
\subsection{A time-dependent question with the same microscopic model}
The transport problem asks how a persistent gradient drives a small departure
from equilibrium. A complementary question is what happens if a departure is
prepared and then allowed to relax, without any spatial gradient or external
force. Around a fixed homogeneous equilibrium, the linearized equation is
\begin{equation}
 w\partial_t\chi=\qtL\chi.
 \label{qt:homogeneous}
\end{equation}
The weight $w$ is time independent in this calculation. Matching has already
removed the mass, momentum and energy modes: the perturbation changes a heat
flux or a shear stress, not the equilibrium parameters.

In either of the two finite trial spaces, write
$\chi_N(t,\qtv)=\sum_nc_n(t)\Phi_n(\qtv)$ and test
Eq.~\eqref{qt:homogeneous} against $\Phi_m$. Define the Gram matrix
\begin{equation}
 H_{mn}=\int w\Phi_m\Phi_n\qtd^3v.
\end{equation}
The time derivative contributes $\sum_nH_{mn}\dot c_n$; the collision term
contributes $-\sum_nM_{mn}c_n$. Therefore
\begin{equation}
 \boxed{H\dot{\boldsymbol c}=-M\boldsymbol c.}
 \label{qt:ode}
\end{equation}
This is an actual evolution generated by the collision operator projected
onto the chosen space. Its rates are not inserted as adjustable decay times.

\subsection{Generalized eigenvalues and the entropy check}
Since both $M$ and $H$ are positive definite in the matched finite space,
solve the symmetric generalized eigenproblem
\begin{equation}
 M\boldsymbol v_r=\nu_rH\boldsymbol v_r,\qquad
 \boldsymbol v_r^{\mathsf T}H\boldsymbol v_s=\delta_{rs},\qquad \nu_r>0.
\end{equation}
The solution is
\begin{equation}
 \boldsymbol c(t)=\sum_r a_r e^{-\nu_rt}\boldsymbol v_r,\qquad
 a_r=\boldsymbol v_r^{\mathsf T}H\boldsymbol c(0).
 \label{qt:decay}
\end{equation}
Thus a general perturbation is a superposition of decaying modes, rather than
necessarily a single exponential. The slowest mode that is actually excited
dominates sufficiently late times.

The quadratic relative-entropy perturbation is
\begin{equation}
 \mathcal E_N(t)=\frac12\int w\chi_N^2\qtd^3v
 =\frac12\boldsymbol c^{\mathsf T}H\boldsymbol c.
\end{equation}
It is the second-order term of the quantum $H$ functional around equilibrium,
up to the convention for conserved additive terms. Differentiating and using
Eq.~\eqref{qt:ode},
\begin{equation}
 \frac{\mathrm d\mathcal E_N}{\mathrm dt}
 =\boldsymbol c^{\mathsf T}H\dot{\boldsymbol c}
 =-\boldsymbol c^{\mathsf T}M\boldsymbol c\leq0.
 \label{qt:entropydecay}
\end{equation}
This is the finite-dimensional counterpart of collision-bracket dissipation,
Eq.~\eqref{qt:bracket}. It gives a direct diagnostic for the time evolution.
The positive eigenvalues demonstrate decay in this trial space; they do not
prove a spectral gap for the full infinite-dimensional quantum operator.

\subsection{The initial perturbation and the observable}
For each sector we choose $\boldsymbol c(0)=H^{-1}\boldsymbol d$ and measure
$J(t)=\boldsymbol d^{\mathsf T}\boldsymbol c(t)$. This is a specified,
reproducible perturbation aligned with the projected transport source.
The normalization of its amplitude is immaterial for the plotted ratios.
For the thermal sector, matching makes the source moment equivalent to the
heat-flux moment; for the shear sector it is the off-diagonal stress moment.
With this choice $a_r=\boldsymbol v_r^{\mathsf T}\boldsymbol d$, so
\begin{equation}
 J(t)=\sum_r a_r^2e^{-\nu_rt},\qquad
 \mathcal E_N(t)=\frac12\sum_r a_r^2e^{-2\nu_rt}.
\end{equation}
In particular both quantities decay monotonically in these runs. Another
initial perturbation need not make every individual flux monotone, even
though the quadratic entropy still decreases.

Integration over time also gives an internal consistency identity,
\begin{equation}
 \int_0^\infty J(t)\qtd t
 =\sum_r\frac{a_r^2}{\nu_r}
 =\boldsymbol d^{\mathsf T}M^{-1}\boldsymbol d.
 \label{qt:timeintegral}
\end{equation}
The last expression is the corresponding static transport coefficient in
our units. This ties stationary response to relaxation in the same discretized
model; it is not an independent molecular-dynamics or Kubo calculation.

\begin{figure}[htbp]\centering\qtfig{relaxation}
\caption{Homogeneous linear relaxation with $N=4$ and the baseline kernel
$B=g/4$. Time is measured in $t_0=(nd^2v_0)^{-1}$. Upper panels show normalized
source moments; lower panels show normalized quadratic relative entropy.
The equilibrium is fixed, and the perturbation evolves according to
$H\dot c=-Mc$.}
\label{qt:relaxfig}\end{figure}
\begin{table}[htbp]\centering
\begin{tabular}{llrr}\toprule
Statistics&$z$&Smallest thermal $\nu_rt_0$&Smallest shear $\nu_rt_0$\\\midrule
\qtinput{\qtroot/data/relaxation_table.tex}
\bottomrule\end{tabular}
\caption{Smallest rates in the $N=4$ matched trial spaces. These finite-space
rates should not be identified with a proven full-operator gap.}
\end{table}
Figure~\ref{qt:relaxfig} is a deterministic simulation of a linearized
distribution. It does not track particles, and it does not solve the nonlinear
kinetic equation for large departures. Polynomial tangent perturbations also
do not automatically ensure global positivity of a reconstructed finite-amplitude
distribution at arbitrarily large velocity. The interpretation here is linear response.

\subsection{A transverse wave: recovering viscosity from dynamics}
We can also retain spatial transport. Consider a small flow in the $x$
direction varying along $y$. A real spatial Fourier mode corresponds to
$u_x(y,t)\propto\operatorname{Re}[\widehat u_x(t)e^{iky}]$.
Its kinetic equation is
\begin{equation}
 w(\partial_t\chi+v_y\partial_y\chi)=\qtL\chi.
\end{equation}
In addition to the $N$ thermal vectors and $N$ shear tensors, include the
conserved function $v_x$. It carries the mean flow and must not be removed
when simulating momentum transport. The projected space has dimension $2N+1$.
Define the streaming matrix
\begin{equation}
 S_{mn}=\int w\Phi_mv_y\Phi_n\qtd^3v.
\end{equation}
For the Fourier amplitudes the equation becomes
\begin{equation}
 \boxed{H\dot{\boldsymbol c}=-(M+ikS)\boldsymbol c.}
 \label{qt:waveode}
\end{equation}
We initialize the conserved velocity component and set the other coefficients
to zero. Streaming then creates a shear stress, and collisions relax it.
The coefficient $\mu$ is not supplied to this kinetic evolution.

At long wavelengths, the hydrodynamic prediction is
\begin{equation}
 \partial_tu_x=\frac\mu\rho\partial_y^2u_x,
 \qquad \widehat u_x(t)\simeq\widehat u_x(0)e^{-(\mu/\rho)k^2t}.
\end{equation}
If $s(k)$ is the slow eigenvalue of $-H^{-1}(M+ikS)$, the dynamic estimate is
\begin{equation}
 \mu_{\mathrm{dyn}}(k)=-\rho\frac{\operatorname{Re}s(k)}{k^2}.
\end{equation}
Agreement with the static coefficient as $k\to0$ checks the coupling between
streaming, collisions and the stress moment. It is a consistency test within
the same projected model, rather than a second independent collision solver.

\begin{figure}[htbp]\centering\qtfig{shear_wave}
\caption{Left: kinetic transverse-mode amplitude at $k=0.2$, compared with
the Navier--Stokes exponential (dotted). A common classical time scale is used
for all curves. Right: viscosity inferred from the slow eigenvalue approaches
the static Galerkin value as $k$ decreases. This transverse mode diffuses; it
is not the oscillatory sound mode of the following section.}
\label{qt:wavefig}\end{figure}
\begin{table}[htbp]\centering
\begin{tabular}{lrrr}\toprule
Statistics&$\mu_4$&$\mu_{\mathrm{dyn}}(0.05)$&Relative difference [\%]\\\midrule
\qtinput{\qtroot/data/wave_table.tex}
\bottomrule\end{tabular}
\caption{Static and dynamic viscosities for the three states in
Fig.~\ref{qt:wavefig}. The largest listed difference is below $0.01\%$.}
\end{table}
At finite $k$, streaming couples to angular sectors absent from this trial
space. The displayed finite-wave-number evolution is consequently a reduced
kinetic model. Resolving the full kinetic crossover would require convergence
in those additional sectors, not only additional radial polynomials.
